% Source: https://artofproblemsolving.com/community/c3h1372981_current_problems
% Stable identity: aops-topic-1372981-q01
Let the boat's speed in still water be $b$ and the current's speed be $c$. With the current the boat moves at $b+c$, against it at $b-c$. In equal times, distances are proportional to speeds, so
\[
b+c=1.5(b-c)\iff 2.5c=0.5b\iff b=5c .
\]
So the boat's speed in still water is $\boxed{5\text{ times the speed of the current}}$. (Without a numerical speed for the current, only this ratio can be found; for example, a current of $2$ mph gives a boat speed of $10$ mph, and then $12=1.5\cdot8$.)
